\emph{Preliminaries.} The blades are a basis of $\Cl_{\Q}(2,2)$ (Theorem \claimref{synthesis.signed-directions-and-the-integer-clifford-algebra}), and $\dagger$ is an anti-automorphism with $\langle x^{\dagger}\rangle_{0} = \langle x\rangle_{0}$.

\emph{(i).} Computed with the product of $\Cl_{\Q}(2,2)$ in the check: the four $A_{i}$, $B'_{j}$ are self-adjoint of square one, each $A_{i}$ commutes with each $B'_{j}$, so $T_{ij}$ is self-adjoint of square one and $E_{ij} = \tfrac12(1 + T_{ij})$ is a projector; $\langle T_{ij}\rangle_{0} = 0$ gives $\langle E_{ij}\rangle_{0} = \tfrac12$. $T_{ij} = 2E_{ij} - 1$ gives the formula for $C$. The commutation table, $D_{k}^{2} = 2 + 2K$, $D_{k}^{3} = 4D_{k}$ and $D_{1}D_{2} \ne D_{2}D_{1}$ are computed. $D_{1} = 2E_{11} - 2E_{22}$ and $D_{2} = 2E_{12} + 2E_{21} - 2$ are combinations of commuting projectors.

\emph{(ii).} $C^{2} = 4 + 4K$ and $C^{3} = 8C$ are computed; $K$ is self-adjoint with $K^{2} = 1$ and $K \ne \pm 1$, so $C^{2} \ne 8$ and $C \ne 0$. Suppose $C = \sum_{i=1}^{n}c_{i}E_{i}$ with rational $c_{i}$ and pairwise commuting projectors $E_{i}$. For $\sigma \in \{0,1\}^{n}$ put $F_{\sigma} = \prod_{i}E_{i}^{(\sigma_{i})}$ with $E^{(1)} = E$, $E^{(0)} = 1 - E$. The $F_{\sigma}$ are idempotents with $F_{\sigma}F_{\sigma'} = 0$ for $\sigma \ne \sigma'$ and $\sum_{\sigma}F_{\sigma} = \prod_{i}(E_{i} + 1 - E_{i}) = 1$, and $CF_{\sigma} = \lambda_{\sigma}F_{\sigma}$ with $\lambda_{\sigma} = \sum_{i}c_{i}\sigma_{i} \in \Q$. So $g(C) = 0$ for $g(t) = \prod(t - \lambda)$, the product over the distinct values $\lambda$ of the $\lambda_{\sigma}$. $C$ also satisfies $t^{3} - 8t = t(t^{2} - 8)$, and $t^{2} - 8$ has no rational root ($8q^{2} = p^{2}$ is impossible for integers, the exponent of $2$ being odd on the left and even on the right). So the greatest common divisor of $g$ and $t^{3} - 8t$ in $\Q[t]$ is $1$ or $t$, and $C$ satisfies it: $1 = 0$ or $C = 0$, both false. For the last sentence: the spectral projectors of the single eigenvalues $0$, $2\sqrt{2}$, $-2\sqrt{2}$ of $C$ in the algebra over the reals are $1 - \tfrac18C^{2}$ and $\tfrac1{16}C^{2} \pm \tfrac{\sqrt{2}}{8}C$ (polynomials in $C$, by $C^{3} = 8C$); the first has rational coefficients and equals $\tfrac12(1 - K)$, the other two do not, because $C$ and $C^{2}$ are linearly independent over $\Q$. A spectral projector of a set of eigenvalues is a sum of these; besides $0$ and $1$ the sums are the three above, $\tfrac18C^{2} = \tfrac12(1 + K)$ (the pair $\pm 2\sqrt{2}$, rational) and $1 - \tfrac1{16}C^{2} \pm \tfrac{\sqrt{2}}{8}C$ (the kernel with one of the two, not rational, for the same reason).

\emph{(iii).} By \claimref{synthesis.additive-weights-on-the-rational-projectors-of-two-modes-are-linear}, clause (iii), with values in $\Q$, $\mu$ is the restriction to the projectors of a $\Q$-linear $L$. So $\sum c_{i}\mu(E_{i}) = \sum c_{i}L(E_{i}) = L(x)$. For self-adjoint $T$ of square one, $\tfrac12(1 \pm T)$ are orthogonal projectors and $\kappa_{\mu}(T) = L(\tfrac12(1 + T)) - L(\tfrac12(1 - T)) = L(T)$; summing, $S(\mu) = L(C)$.

\emph{(iv).} By clause (iv) of the same theorem a record weight is the restriction of a state: $L(1) = 1$ and $L(xx^{\dagger}) \ge 0$. For rational $\tau$, $C - \tau$ is self-adjoint, so $0 \le L((C - \tau)(C - \tau)^{\dagger}) = L(C^{2}) - 2\tau L(C) + \tau^{2}$; with $\tau = L(C)$, $L(C)^{2} \le L(C^{2}) = 4 + 4L(K)$. Since $K^{\dagger} = K$ and $K^{2} = 1$, $(1 - K)(1 - K)^{\dagger} = 2(1 - K)$, so $L(1 - K) \ge 0$ and $L(K) \le 1$. Hence $S(\mu)^{2} \le 8$, and $S(\mu)$ is rational while $8$ is not the square of a rational. The value $\tfrac{14}{5}$ is that of the state $x \mapsto 4\langle\Pi x\rangle_{0}$ for the atom $\Pi$ of the ray $(-1, 2, -2, -1)$ in the notation of the theorem for two modes, with correlators $\tfrac35, \tfrac45, \tfrac45, -\tfrac35$ (computed); the check also verifies $S = L(C)$, $S^{2} \le 4 + 4L(K) \le 8$ and $S^{2} < 8$ for the states of all atoms with integer entries of absolute value at most $3$, among which $\tfrac{14}{5}$ is the largest value.

\emph{(v).} $\lambda$ is linear, so it adds on orthogonal projectors, and $\lambda(1) = 1 + \langle C\rangle_{0} = 1$. The sixteen values, the correlators $1, 1, 1, -1$ and the value $-\tfrac14$ at the atom of the ray $(1, 0, 0, 1)$ are computed in the check (the last is clause (iv) of the theorem for two modes).
